\chapter{Progettazione degli oggetti}
\label{cap:sdd-object-design}

Questo capitolo scende dal livello dei sottosistemi a quello delle classi.
Fissa le convenzioni con cui il codice è scritto, descrive l'organizzazione
in package, documenta i contratti delle classi principali e discute i pattern
adottati — e quelli valutati e scartati.

\section{Compromessi di progettazione degli oggetti}

\begin{longtable}{@{}L{3.6cm} L{4.2cm} L{6.4cm}@{}}
\toprule
\textbf{Compromesso} & \textbf{Alternativa scartata} & \textbf{Motivazione} \\
\midrule
\endfirsthead
\toprule
\textbf{Compromesso} & \textbf{Alternativa scartata} & \textbf{Motivazione} \\
\midrule
\endhead
\bottomrule
\endlastfoot

Entità come classi con accessori &
Strutture dati anonime, senza comportamento &
Le classi permettono di collocare nel dominio le poche operazioni che gli
appartengono davvero, come «la risposta corretta di questa domanda». Il
prezzo è più codice di accesso. \\

Un DAO per aggregato &
Un DAO per tabella &
Domande e risposte non hanno ciclo di vita autonomo: separarle in DAO propri
produrrebbe classi che nessuno può usare correttamente da sole e
moltiplicherebbe le interrogazioni. \\

Iniezione tramite parametri con valore predefinito &
Un contenitore di dipendenze &
Il valore predefinito rende testabile ogni classe senza introdurre un
framework e senza che il codice di produzione debba dichiarare un grafo di
dipendenze. \\

Eccezioni tipizzate &
Valori di ritorno che descrivono l'esito &
Le eccezioni tolgono dalle rotte la propagazione manuale dell'errore, al
prezzo di un percorso di controllo non visibile nella firma dei metodi. \\

\end{longtable}

\section{Linee guida di documentazione delle interfacce}
\label{sec:sdd-linee-guida}

Le convenzioni seguenti valgono per tutto il codice del server e sono
verificate, per quanto possibile, dall'analisi statica.

\begin{description}[leftmargin=1.4cm,style=nextline]
  \item[Lingua] I concetti del dominio conservano il nome italiano
    (\id{Utente}, \id{Tutorial}, \id{Svolgimento}); i termini tecnici restano
    in inglese (\id{Service}, \id{Dao}, \id{middleware}, \id{findById}).
    La regola tiene allineati codice e glossario senza tradurre termini che
    non hanno un equivalente italiano consolidato.
  \item[Nomi] Classi in \id{PascalCase}, metodi e variabili in
    \id{camelCase}, costanti in \id{MAIUSCOLO\_CON\_UNDERSCORE}, tabelle e
    colonne in \id{snake\_case}. I metodi dei DAO usano i verbi
    \id{find}, \id{create}, \id{update}, \id{delete}; quelli dei servizi
    usano il verbo del dominio (\id{creaQuiz}, \id{eseguiQuiz}).
  \item[Visibilità] I campi delle entità sono privati e accessibili tramite
    metodi; è pubblico solo ciò che fa parte del contratto.
  \item[Documentazione] Ogni classe e ogni metodo pubblico hanno un commento
    che ne dichiara lo scopo, i parametri, il valore restituito e le
    eccezioni sollevate. I commenti spiegano \emph{perché}, non
    \emph{che cosa}: ciò che il codice già dice non viene ripetuto.
  \item[Tipi] Nessun uso di \id{any}; il controllo dei tipi è attivo in
    modalità stretta e le variabili non utilizzate sono errori.
\end{description}

\section{Organizzazione in package}

\diagramma{0.62}{pkg-packages}{Package del server e direzione delle
dipendenze.}{fig:sdd-pkg}

La struttura è la stessa per tutti i sottosistemi. Le entità sono raggruppate
per sottosistema, in sottocartelle di \file{entity}; rotte, servizi e DAO sono
invece raggruppati per livello. In questo modo la decomposizione per
sottosistema resta leggibile dai nomi dei file, e la struttura a livelli resta
visibile nell'albero delle cartelle.

Il grafo delle dipendenze è aciclico e orientato verso il basso. Il package
\file{entity} non dipende da nulla: le entità possono essere costruite e
verificate senza alcuna infrastruttura.

\section{Contratti delle classi principali}

\subsection{Livello di persistenza}

\diagramma{0.98}{cd-design-quiz}{Progettazione degli oggetti del sottosistema
Gestione quiz.}{fig:sdd-cd-quiz}

\id{QuizDao} è l'esempio più significativo: gestisce l'intero aggregato
quiz-domande-risposte. La lettura avviene con un'unica interrogazione che
unisce le tre tabelle e ricostruisce il grafo di oggetti in memoria; la
scrittura accetta una connessione transazionale, perché coinvolge più
tabelle e deve essere atomica.

\begin{longtable}{@{}L{5.6cm} L{8.8cm}@{}}
\toprule
\textbf{Operazione} & \textbf{Contratto} \\
\midrule
\endfirsthead
\toprule
\textbf{Operazione} & \textbf{Contratto} \\
\midrule
\endhead
\bottomrule
\endlastfoot

\id{findById(id): Quiz | null} &
Restituisce l'aggregato completo o \id{null}. Una sola interrogazione. \\

\id{findByTutorial(tutorialId): Quiz | null} & Come sopra, per tutorial. \\

\id{create(quiz, connection): number} &
Inserisce quiz, domande e risposte assegnando gli identificativi alle entità
passate. \emph{Richiede} una connessione transazionale. \\

\id{replaceDomande(quizId, domande, connection)} &
Elimina le domande esistenti — le risposte cadono in cascata — e inserisce
le nuove. \emph{Richiede} una connessione transazionale. \\

\id{delete(id): boolean} &
Restituisce \id{false} se il quiz non esisteva. \\

\end{longtable}

\subsection{Serializzazione}
\label{sec:sdd-dto}

Il package \file{dto} contiene funzioni pure che traducono un'entità nella
rappresentazione destinata al client. Nessuna rotta restituisce
direttamente un'entità.

Questa regola non è una preferenza stilistica: è il meccanismo con cui due
requisiti di sicurezza vengono fatti rispettare in un punto solo
(\id{DG\_1}, \id{DG\_2}).

\begin{longtable}{@{}L{4.0cm} L{10.4cm}@{}}
\toprule
\textbf{Funzione} & \textbf{Garanzia} \\
\midrule
\endfirsthead
\toprule
\textbf{Funzione} & \textbf{Garanzia} \\
\midrule
\endhead
\bottomrule
\endlastfoot

\id{toUtenteDTO} &
Omette l'hash della password. È l'unico modo in cui un utente lascia il
server, quindi la credenziale non può sfuggire per dimenticanza in una rotta
(\id{RNF\_F\_1}). \\

\id{toQuizDTO} &
Omette il campo che distingue la risposta corretta. Rende
strutturalmente impossibile erogare un quiz con le soluzioni
(\id{RNF\_F\_6}). \\

\id{toTutorialDTO}, \id{toFeedbackDTO}, \id{toObiettivoDTO}, \id{toBadgeDTO} &
Convertono le date in formato ISO 8601 e normalizzano i valori numerici
provenienti dal DBMS. \\

\end{longtable}

\subsection{Gestione degli errori}
\label{sec:sdd-errori}

\diagramma{0.72}{cd-errori}{Gerarchia degli errori applicativi.}{fig:sdd-errori}

I servizi non restituiscono descrizioni dell'esito: sollevano eccezioni che
portano con sé lo stato HTTP e un codice stabile. Un unico middleware le
traduce in risposte. La conseguenza pratica è che una rotta si riduce a tre
righe — estrarre i parametri, invocare il servizio, serializzare — e che il
percorso d'errore è identico in tutto il sistema.

Un errore che non discende da \id{AppError} è considerato imprevisto:
produce \id{500} con un messaggio generico, mentre il dettaglio resta nei
log del server (\id{RNF\_A\_2}).

\subsection{Iniezione delle dipendenze}
\label{sec:sdd-iniezione}

Ogni servizio riceve i propri collaboratori come parametri del costruttore,
con un valore predefinito che costruisce l'implementazione reale:

\begin{lstlisting}[language=Java,caption={Iniezione con valore predefinito.},
                   label={lst:sdd-di}]
export class QuizService {
  constructor(
    private readonly quizDao: QuizDao = new QuizDao(),
    private readonly svolgimentoDao: SvolgimentoDao = new SvolgimentoDao(),
    private readonly utenteDao: UtenteDao = new UtenteDao(),
    private readonly tutorialDao: TutorialDao = new TutorialDao(),
    private readonly obiettivoService: ObiettivoService = new ObiettivoService(),
  ) {}
  ...
}
\end{lstlisting}

In esercizio la classe si costruisce da sola; nei test riceve dei doppi.
Nessun contenitore di dipendenze, nessuna configurazione, e soprattutto
nessun bisogno per i test di raggiungere campi privati per sostituire un
collaboratore (\id{DG\_3}).

Alla radice dell'applicazione la funzione \id{creaApp} riceve l'insieme dei
servizi: è il punto in cui il grafo viene composto, ed è anche il punto in
cui i test di integrazione lo sostituiscono per intero.

\section{Pattern di progettazione adottati}
\label{sec:sdd-pattern}

\subsection{Data Access Object}

\textbf{Problema.} Il codice di dominio non deve conoscere SQL, e il livello
di persistenza deve poter essere sostituito nei test.

\textbf{Soluzione.} Ogni aggregato ha un DAO che espone operazioni in termini
di entità e nasconde le interrogazioni. I servizi dipendono dall'interfaccia
del DAO, non dalla sua realizzazione.

\textbf{Conseguenze.} L'intera suite di test gira senza database
(\id{DG\_3}). In cambio, le interrogazioni non sono verificate dai test
automatici: è il compromesso già dichiarato in \cref{sec:sdd-tradeoff}.

\subsection{Facade (lato client)}

\textbf{Problema.} Le pagine dell'interfaccia non devono conoscere URL,
formato del trasporto e gestione degli errori di rete: se lo facessero, un
cambiamento nell'API si propagherebbe in decine di componenti.

\textbf{Soluzione.} Il modulo \file{src/api} espone un unico oggetto
organizzato per sottosistema. Le pagine invocano
\id{api.quiz.consegna(\dots)} e ricevono dati già tipizzati o un errore
già normalizzato~\cite{gamma1994}.

\textbf{Conseguenze.} Il confine fra client e server è attraversato in un
solo punto. Il client non importa alcun modulo del server: i tipi di
trasporto sono ridichiarati in \file{src/types}, duplicazione accettata
consapevolmente in cambio dell'indipendenza fra le due applicazioni.

\subsection{Chain of Responsibility (middleware)}

\textbf{Problema.} Autenticazione, autorizzazione, validazione e gestione
degli errori riguardano tutte le rotte e non appartengono a nessun
sottosistema.

\textbf{Soluzione.} Sono realizzate come middleware componibili, attraversati
in sequenza: ciascuno può concludere la richiesta o passarla al successivo.
Una rotta dichiara le proprie esigenze componendoli
(\id{requireAdmin}, poi i validatori, poi il gestore).

\textbf{Conseguenze.} Il controllo degli accessi è dichiarativo e leggibile
accanto alla rotta; la matrice di \cref{sec:sdd-accessi} è verificabile
leggendo i router. Il rovescio è che dimenticare un middleware su una rotta
la lascia scoperta: per questo la matrice è coperta da test dedicati.

\subsection{Template Method (transazioni)}

\textbf{Problema.} Aprire una transazione, effettuare commit, gestire il
rollback e rilasciare la connessione è una sequenza invariante che sarebbe
ripetuta e sbagliata a ogni occorrenza.

\textbf{Soluzione.} La funzione \id{withTransaction} fissa quella sequenza e
riceve come parametro il lavoro variabile da eseguire.

\textbf{Conseguenze.} È impossibile dimenticare un rollback o non rilasciare
una connessione. La responsabilità che resta al chiamante è propagare la
connessione ai DAO che invoca.

\section{Pattern valutati e non adottati}

Documentare i pattern scartati è utile quanto documentare quelli adottati:
evita che la loro assenza sia scambiata per una svista.

\begin{longtable}{@{}L{2.8cm} L{5.4cm} L{6.2cm}@{}}
\toprule
\textbf{Pattern} & \textbf{Applicazione ipotizzata} & \textbf{Perché è stato scartato} \\
\midrule
\endfirsthead
\toprule
\textbf{Pattern} & \textbf{Applicazione ipotizzata} & \textbf{Perché è stato scartato} \\
\midrule
\endhead
\bottomrule
\endlastfoot

Proxy &
Caricare i tutorial in forma ridotta nel catalogo e recuperare il contenuto
completo solo all'apertura. &
Il catalogo contiene pochi elementi e la selezione delle colonne risolve lo
stesso problema con una clausola SQL. Un Proxy aggiungerebbe una classe e un
livello di indirezione per un beneficio che il database offre già. \\

Observer &
Notificare l'assegnazione di un badge a più destinatari interessati. &
Esiste un solo destinatario, la risposta HTTP in corso. Il pattern avrebbe
senso con notifiche via email o eventi persistenti, che sono fuori
perimetro. \\

Strategy &
Rendere sostituibile il criterio di superamento del quiz. &
Il criterio è unico e fissato dal dominio. La soglia è già una costante
condivisa: renderla polimorfa non risponderebbe ad alcuna necessità reale. \\

Repository con Unit of Work &
Tracciare le modifiche alle entità e persisterle in blocco. &
Le transazioni del sistema coinvolgono poche entità note a priori.
\id{withTransaction} copre il caso d'uso senza introdurre un livello di
mappatura degli oggetti. \\

Singleton &
Garantire un'unica istanza dei servizi. &
I servizi sono privi di stato: istanziarli più volte non ha conseguenze, e
il Singleton renderebbe più difficile sostituirli nei test. Il solo elemento
realmente unico, il pool di connessioni, è già un modulo con inizializzazione
differita. \\

\end{longtable}

\begin{nota}
Un pattern citato in un documento di progettazione ma assente dal codice non
è un dettaglio redazionale: è una descrizione errata del sistema, e induce
chi legge a cercare una struttura che non esiste. Per questo i pattern di
\cref{sec:sdd-pattern} corrispondono a classi e funzioni effettivamente
presenti, mentre quelli valutati e non adottati sono dichiarati come tali
invece di essere taciuti.
\end{nota}
